\chapter{طراحی و پیاده‌سازی تجمیع میانگین وزنیِ احرازاصالت‌شده با \en{SBHE-Benaloh}}
\label{chap:sbhe}
\section{هدف و مرز پروتکل}
در یک سامانۀ چندکلاینتی، علاوه بر محاسبۀ مجموع وزن‌دار، باید مشخص باشد هر مشارکت از کدام هویت مجاز آمده، در مسیر تغییر نکرده و متعلق به دور جاری است. خاصیت همریختی به‌تنهایی این نیازها را پاسخ نمی‌دهد. ساختار \en{SBHE-Benaloh} در این مخزن، متن رمز بنالوح را با امضای \en{Ed25519} روی یک محمولۀ مشخص همراه می‌کند. امضا پیش از ورود متن رمز به تجمیع بررسی می‌شود و در عملگر جبری متن رمز ضرب نمی‌شود.

این طرح در سطح کد، یک پوشش احراز اصالت برای تجمیع همریخت است؛ ادعای یک سازوکار استانداردِ مستقل \en{signcryption} با اثبات امنیت ترکیبی در یک مدل رسمی مطرح نمی‌شود. همچنین سرور در پیاده‌سازی حاضر مالک کلید خصوصی بنالوح است. او طبق مسیر برنامه فقط مجموع را رمزگشایی می‌کند، اما از نظر رمزنگاری قادر به رمزگشایی یک متن رمز منفرد نیز هست. در نتیجه، محرمانگی در برابر شنودگرِ فاقد کلید با محرمانگی در برابر سرورِ دارندۀ کلید یکسان نیست.

\begin{figure}[htbp]
\centering\begin{tikzpicture}[
  node distance=5mm,
  stage/.style={draw=blue!50!black,fill=blue!3,rounded corners=3pt,
    inner sep=6pt,font=\small},
  flow/.style={-{Latex[length=2.2mm]},thick,draw=blue!50!black}]
\node[stage] (encode) {\diagramtext{کلاینت مجاز: بررسی مقدار و نمایش دقیق با مقیاس $S=10000$\\
  $m_i=Sx_i,\quad 0\le m_i\le S,\quad w_i\in\Z_{>0}$}};
\node[stage,below=of encode] (encrypt) {\diagramtext{رمزگذاری بنالوح با تصادف تازه\\
  $C_i=y^{m_i}u_i^r\bmod n,\quad u_i\in\Z_n^*$}};
\node[stage,below=of encrypt] (sign) {\diagramtext{ساخت پیام قطعی شامل شناسۀ طرح، کلاینت، دور، \en{nonce}، متن رمز و وزن\\
  $M_i=\operatorname{Serialize}(\mathrm{scheme},ID_i,RID,\nu_i,C_i,w_i)$\\
  $\sigma_i=\Sign_{SK_i}(M_i)$؛ امضای مستقیم با \en{Ed25519}}};
\node[stage,below=of sign] (verify) {\diagramtext{سرور: بررسی هویت مجاز، دور، وزن، بازپخش و عضویت گروهی\\
  $\Verify_{PK_i}(M_i,\sigma_i)=\mathrm{true}$\\
  رد محمولۀ نامعتبر پیش از تجمیع؛ ثبت مشارکت‌های پذیرفته‌شده}};
\node[stage,below=of verify] (aggregate) {\diagramtext{تجمیع مؤلفۀ رمز‌شدۀ کلاینت‌های پذیرفته‌شده\\
  $W=\sum_{i\in\mathcal A}w_i,\quad SW<r,\quad
  C_{\mathrm{agg}}=\prod_{i\in\mathcal A}C_i^{w_i}\bmod n$}};
\node[stage,below=of aggregate] (decrypt) {\diagramtext{رمزگشایی و بازیابی میانگین\\
  $M=\Dec_B(C_{\mathrm{agg}}),\quad \bar x=M/(SW)$}};
\node[stage,below=of decrypt] (limit) {\diagramtext{سرورِ مالک کلید خصوصی توان رمزگشایی پیام منفرد را نیز دارد\\
  امضا، صحت مقدار ادعاییِ یک کلاینت مجاز را اثبات نمی‌کند}};
\foreach \a/\b in {encode/encrypt,encrypt/sign,sign/verify,verify/aggregate,aggregate/decrypt,decrypt/limit}
  \draw[flow] (\a) -- (\b);
\end{tikzpicture}

\caption{جریان اجرای \en{SBHE-Benaloh}: امضای مستقیم پیام بایتی و بررسی محموله پیش از تجمیع مؤلفۀ رمز‌شده}
\label{fig:sbhe-workflow}
\end{figure}

\section{مدل سیستم و نقش اعضا}
سرور یک کلید بنالوح تولید و کلید عمومی $(n,y,r)$ را منتشر می‌کند. هر کلاینت $i$ یک شناسۀ یکتا $ID_i$ و جفت کلید امضای $(SK_i,PK_i)$ دارد. پیش از مشارکت، ارتباط معتبر میان شناسه و کلید عمومی امضا باید در فهرست مجاز سرور ثبت شود. جایگزینی کلید یک هویت مجاز توسط مهاجم، این پیش‌فرض را نقض می‌کند و با بررسی جبری امضا قابل جبران نیست.

هر پیام به شناسۀ دور $RID$ و یک مقدار یک‌بارمصرف $\nu_i$ متصل است. مقدار ورودی کلاینت $x_i\in[0,1]$ و وزن آن $w_i\in\Z_{>0}$ است. در آزمایش‌های کوچک، مقدارها دادۀ عددی آزمایشی هستند. برای کاربرد مشارکتی، $x_i$ می‌تواند نسخۀ کوانتیزه‌شدۀ احتمال خروجی مدل باشد و $w_i$ تعداد نمونۀ محلی یا وزن صحیحِ تعریف‌شدۀ دیگری باشد؛ صحت این تفسیر به سیاست سامانه وابسته است.

فهرست اعضای مجاز، سیاست تعیین دور، نگهداری کلید خصوصی و پایداری دفتر \en{nonce} جزو مدل عملیاتی‌اند. این مؤلفه‌ها باید جدا از الگوریتم بنالوح تعریف شوند. در نسخۀ فعلی، فهرست کلیدها و مجموعۀ \en{nonce}های مصرف‌شده در حافظۀ شیء سرور نگهداری می‌شوند؛ بازیابی پس از راه‌اندازی مجدد به ذخیره‌سازی پایدار و سیاست دور نیاز دارد.

\section{بازنمایی عددی و مدیریت دقت}
\label{sec:fixed-point}
بنابراین که فضای پیام بنالوح صحیح است، ضریب مقیاس $S=10^4$ تعریف می‌شود. در مسیر فعلی کد، $x_i$ باید حداکثر چهار رقم اعشار داشته باشد تا
\begin{equation}
m_i=Sx_i\in\{0,\ldots,S\}
\label{eq:fixed-point}
\end{equation}
بدون گرد کردن پنهان محاسبه شود. برای نمونه، $0.7314$ به عدد صحیح $7314$ تبدیل می‌شود. تابع \en{encode\_value} ورودیِ دارای اعشار بیشتر، مانند $0.12345$، و مقدار خارج از بازه را رد می‌کند. استفاده از \en{Decimal}، هدف نمایش و تقسیم اعشاری را روشن می‌کند و جایگزین کردن تصادفیِ قالب اعشاری با عدد دودویی را کاهش می‌دهد.

احتمال خروجی طبقه‌بند لزوماً چهاررقم‌اعشاری نیست. بنابراین اتصال مستقیم نوت‌بوک به این تابع، بدون تعریف مرحلۀ تبدیل، معتبر نیست. اگر سیاستی صریح برای گرد کردن به نزدیک‌ترین مقدار چهاررقمی تعریف شود و $q_i$ خروجی آن باشد، آنگاه $|q_i-x_i|\le1/(2S)$ و برای وزن‌های مثبت داریم:
\begin{equation}
\left|\frac{\sum_iw_iq_i}{W}-\frac{\sum_iw_ix_i}{W}\right|
\le\frac{1}{2S},\qquad W=\sum_iw_i.
\label{eq:quantization-bound}
\end{equation}
این کران مربوط به میانگین است. اگر احتمال تجمیعی به آستانۀ طبقه‌بندی نزدیک باشد، حتی همین خطای کوچک می‌تواند برچسب را تغییر دهد. در این پایان‌نامه، دقت‌های مشارکتیِ فصل قبل حاصل محاسبۀ آشکارِ اعشاری‌اند؛ برابری برچسب‌ها پس از این تبدیل اندازه‌گیری نشده است.

\section{بن‌مایۀ جبری بنالوح}
کلاینت برای تصادف تازۀ $u_i\in\Z_n^*$ محاسبه می‌کند:
\begin{equation}
C_i=y^{m_i}u_i^r\pmod n.
\end{equation}
ضرب متن‌های رمز متناظر با جمع پیام‌ها و توان‌رسانی به وزن آشکار متناظر با ضرب پیام در وزن است. برای کلاینت‌های پذیرفته‌شدۀ $\mathcal A$، نتیجۀ تجمیع برابر است با
\begin{equation}
C_{\mathrm{agg}}=\prod_{i\in\mathcal A}C_i^{w_i}\pmod n.
\label{eq:sbhe-aggregate}
\end{equation}
امضای دیجیتال هیچ تغییری در این رابطه ایجاد نمی‌کند. پارامتر $r$ باید با بیشینۀ مجموع صحیح سازگار باشد؛ اندازۀ پیمانۀ رمز $n$ و اندازۀ فضای پیام $r$ دو پارامتر با نقش متفاوت‌اند.

\section{ساختار محموله و پیام امضاشده}
محمولۀ کلاینت به‌صورت زیر تعریف می‌شود:
\begin{equation}
E_i=(ID_i,RID,\nu_i,C_i,w_i,\sigma_i).
\label{eq:sbhe-envelope}
\end{equation}
پیام بایتی $M_i$ از نمایش قطعیِ شناسۀ طرح و مؤلفه‌های محموله ساخته می‌شود:
\begin{equation}
M_i=\operatorname{Serialize}(\mathrm{scheme},ID_i,RID,\nu_i,C_i,w_i),
\qquad \sigma_i=\Sign_{SK_i}(M_i).
\label{eq:sbhe-signature}
\end{equation}
در کد، \en{Serialize} با \en{JSON} دارای کلیدهای مرتب‌شده، جداکننده‌های بدون فاصله، تبدیل متن رمز به رشتۀ ده‌دهی و کدگذاری \en{UTF-8} انجام می‌شود. شناسۀ نسخۀ طرح \en{SBHE-Benaloh-Ed25519-v1} نیز در پیام امضاشده حضور دارد. این صورت دقیق با نوشتن مبهمِ چسباندن رشته‌ها تفاوت دارد؛ مرز فیلدها باید برای فرستنده و گیرنده یکسان باشد.

پیاده‌سازی، بایت‌های همین پیام را مستقیماً به \en{Ed25519.sign} می‌دهد. رابطۀ اجراشده، امضای دستیِ $H(M_i)$ یا گونۀ \en{Ed25519ph} نیست؛ هش داخلیِ الگوریتم \en{Ed25519} در تعریف خود الگوریتم انجام می‌شود \cite{ed25519}.

وزن باید داخل پیام امضاشده باشد؛ در غیر این صورت مهاجم می‌تواند با نگه داشتن متن رمز و امضای آن، وزن را عوض کند. شناسۀ دور، پیام را به یک اجرا متصل می‌کند و \en{nonce} امکان تشخیص استفادۀ مجدد همان مشارکت را می‌دهد. امضای تمام این فیلدها، تغییر برون‌زاد را قابل تشخیص می‌کند، اما ادعای کلاینت دربارۀ درست بودن وزن را به‌تنهایی اثبات نمی‌کند.

\section{گام‌های کلاینت و سرور}
\subsection{فرایند کلاینت}
کلاینت ابتدا مقدار و وزن را بررسی، مقدار مجاز را با مقیاس $S$ به صحیح تبدیل و آن را با کلید عمومی معتبر بنالوح رمز می‌کند. سپس \en{nonce} تازه می‌سازد، پیام قطعی را از فیلدهای مشخص تشکیل می‌دهد و با کلید خصوصی خود امضا می‌کند. در پایان محمولۀ کامل را می‌فرستد. تصادف رمزگذاری و \en{nonce} دو نقش متفاوت دارند و یکی جایگزین دیگری نیست.

\subsection{پذیرش محموله در سرور}
سرور باید پیش از تجمیع، موارد زیر را بررسی کند:
\begin{enumerate}
\item شناسۀ دورِ محموله با دور مورد انتظار برابر باشد.
\item شناسۀ کلاینت در فهرست مجاز وجود داشته باشد.
\item وزن مطابق قالب، صحیح و مثبت باشد؛ سیاست کاربرد دربارۀ مقدار مجاز وزن نیز جداگانه اعمال شود.
\item سه‌تایی $(ID_i,RID,\nu_i)$ قبلاً مصرف نشده باشد.
\item امضا با کلید ثبت‌شده و همان بایت‌های قطعیِ پیام معتبر باشد.
\item متن رمز در دامنۀ گروه قرار گیرد: $0<C_i<n$ و $\gcd(C_i,n)=1$.
\end{enumerate}
محمولۀ نامعتبر نباید در حاصل‌ضرب وارد شود. تابع \en{accept\_and\_aggregate} همچنین مشارکت تکراری یک شناسه در همان فراخوانیِ دسته‌ای را رد می‌کند. این کنترل درون یک دسته با سیاست «فقط یک مشارکت از هر کلاینت در تمام عمر یک دور» یکسان نیست؛ برای سیاست دوم باید وضعیت دور در کل سامانه نگهداری شود.

ثبت \en{nonce}های مصرف‌شده در کد پس از اعتبارسنجی دسته و ساخت تجمیع انجام می‌شود. این ترتیب از مصرف شدن بخشی از شناسه‌ها در یک دستۀ ناموفق جلوگیری می‌کند. در یک سرور هم‌زمانِ واقعی، بررسی و ثبت باید در برابر رقابت دو درخواست نیز اتمی باشد؛ ساختار درون‌حافظه‌ای یک نمونۀ پژوهشی، به‌تنهایی این نیاز استقرار را حل نمی‌کند.

\subsection{رمزگشایی و شرط ظرفیت پیام}
برای مجموع وزن $W>0$ و مجموع صحیح $M=\sum_{i\in\mathcal A}w_im_i$، پس از رمزگشایی داریم:
\begin{equation}
\bar x=\frac{M}{SW}.
\label{eq:sbhe-decode}
\end{equation}
چون فضای پیام پیمانه‌ای است، شرط کافیِ تضمین نبود سرریز برای $0\le m_i\le S$ چنین است:
\begin{equation}
SW<r.
\label{eq:sbhe-capacity}
\end{equation}
با $r=10000019$ و $S=10000$، مجموع وزن صحیح حداکثر ۱۰۰۰ مجاز است. شرط فوق در مسیر تجمیع کد کنترل می‌شود و عبور از ظرفیت باید با خطا متوقف شود. کوچک بودن تعداد کلاینت‌ها به‌تنهایی کافی نیست؛ مجموع وزن‌ها تعیین‌کننده است. برای وزن‌های تعداد نمونۀ آزمایش مشارکتی، ظرفیت لازم در رابطۀ~\eqref{eq:fl-benaloh-capacity} بسیار بزرگ‌تر است.

\section{اثبات درستی تجمیع}
فرض می‌شود همۀ ورودی‌های پذیرفته‌شده پیام‌های معتبر در بازۀ تعریف‌شده‌اند و شرط~\eqref{eq:sbhe-capacity} برقرار است. با جایگذاری متن رمز کلاینت در رابطۀ تجمیع:
\begin{align}
C_{\mathrm{agg}}
&=\prod_{i\in\mathcal A}(y^{m_i}u_i^r)^{w_i}\pmod n\\
&=y^{\sum_iw_im_i}\left(\prod_iu_i^{w_i}\right)^r\pmod n\\
&=\Enc_{B}\!\left(\sum_iw_im_i\right).
\end{align}
در نتیجه، رمزگشایی مجموع صحیح را در پیمانۀ $r$ به دست می‌دهد. از آنجا که $0\le M\le SW<r$، نتیجۀ پیمانه‌ای همان مجموع صحیحِ مورد نظر است. تقسیم بر $SW$ می‌دهد:
\begin{equation}
\frac{\Dec_B(C_{\mathrm{agg}})}{SW}
=\frac{\sum_iw_i(Sx_i)}{SW}
=\frac{\sum_iw_ix_i}{\sum_iw_i}.
\end{equation}
این اثبات، درستی جبری را تحت پیش‌شرط‌های ورودی نشان می‌دهد؛ نشان نمی‌دهد کلاینت مخرب الزاماً همان احتمالِ حاصل از مدل خود را رمز کرده است. برای این ادعای قوی‌تر، گواه محاسبه یا کنترل اعتماد دیگری لازم است.

\section{تحلیل امنیتی و محدودیت‌های دقیق}
\subsection{هویت، تمامیت و بازپخش}
با فرض محرمانه ماندن کلید امضا و ثبت معتبر کلید عمومی، مهاجم بیرونی نباید بتواند برای شناسۀ مجاز محمولۀ تازۀ معتبر بسازد. تغییر متن رمز، وزن، شناسه، دور یا \en{nonce} بایت‌های پیام را تغییر می‌دهد و باید بررسی امضا را ناموفق کند. کنترل دور و دفتر مصرف \en{nonce} نیز از پذیرش دوبارۀ همان پیام جلوگیری می‌کنند. پایداری و هماهنگی آن دفتر در استقرار واقعی جزو پیش‌فرض‌هاست.

\subsection{مرز محرمانگی}
متن رمز، مقدار منفرد را از طرف فاقد کلید خصوصی پنهان می‌کند، در چارچوب فرض امنیتی و انتخاب پارامتر معتبر. با این حال، مالک کلید خصوصی توان رمزگشایی متن منفرد را دارد. مشاهدۀ میانگین نهایی نیز ممکن است با دانش جانبی یا در گروه‌های بسیار کوچک اطلاعاتی دربارۀ مشارکت‌کننده آشکار کند. بنابراین این ساختار نه حریم خصوصی تفاضلی است و نه یک طرح رمزگشایی آستانه‌ایِ اجراشده.

\subsection{کلاینت مجاز اما مخرب}
امضای معتبر فقط نشان می‌دهد محموله با کلید ثبت‌شده تولید شده است. کلاینت مجاز می‌تواند مقدار یا وزن نادرست را امضا کند. کنترل عضویت $C_i\in\Z_n^*$ ثابت نمی‌کند پیام رمز‌شده در بازۀ $[0,S]$ قرار دارد. کنترل بازه در تابع محلی، در برابر کلاینتی که آن تابع را دور بزند کافی نیست. اثبات بازه، اتصال وزن به شمارش دادۀ تأییدشده، مقاومت در برابر مسموم‌سازی مدل و صحت استنتاج، در نسخۀ فعلی پیاده‌سازی نشده‌اند.

\subsection{هزینه، مدیریت کلید و استقرار}
امضاهای \en{Ed25519} در این پیاده‌سازی جداگانه بررسی می‌شوند و با ضرب متن‌های رمز تجمیع نمی‌شوند. هزینۀ اعتبارسنجی با تعداد مشارکت‌ها رشد می‌کند. افشای کلید خصوصی یک عضو، جایگزینی کلید در رجیستری، ضعف تصادف، نشت کانال جانبی و حذف وضعیت بازپخش، هر یک پیش‌فرضی جداگانه را نقض می‌کنند. اندازه‌های کوچک کلید در آزمایش زمان، پارامتر امنیتی تولیدی نیستند. برای کاربرد عملی باید از کتابخانه‌ها و پارامترهای ممیزی‌شده و مدل تهدید مستقل استفاده شود.

\section{چرا امضا در متن رمز ضرب نمی‌شود؟}
یک ترکیب ساده‌انگارانه ممکن است $C'=C\sigma\pmod n$ را پیشنهاد کند. امضای \en{Ed25519} عضو دارای همان نقش جبری در فضای پیام بنالوح نیست؛ ضرب چنین عاملی می‌تواند رابطۀ رمزگشایی را خراب کند و هیچ قاعدۀ معتبر احراز اصالتی نیز ایجاد نمی‌کند. در ساختار حاضر، محموله یک زوج صرفاً جبریِ دو عضو هم‌نوع نیست؛ مجموعه‌ای از فیلدهای دارای معناست. عملیات همریخت فقط بر مؤلفۀ $C_i$ انجام می‌شود و $\sigma_i$ پیش‌شرط پذیرش همان مؤلفه و وزن آن است.

این جداسازی امکان استفاده از همان قانون جمعی بنالوح را حفظ می‌کند، اما به‌تنهایی اثبات امنیت ترکیبیِ همۀ لایه‌ها نیست. ثبت هویت، قالب قطعی، جداسازی دامنۀ نسخه، انتخاب دور و وضعیت بازپخش، بخش‌های ضروری تحلیل کامل‌اند.

\section{ارزیابی عملکردِ ذخیره‌شده}
\label{sec:sbhe-performance}
منبع این بخش فایل‌های تفصیلی و خلاصۀ \en{Benaloh\_signcryption/results} با تاریخ ثبت‌شدۀ \en{20260823\_180220} است. زمان‌های موجود، آمار اجرای ذخیره‌شده‌اند و به‌عنوان اجرای مجدد روی محیط فعلی معرفی نمی‌شوند. کد \en{demo.py} به‌صورت پیش‌فرض ۹ پیکربندیِ سه اندازۀ کلید و سه تعداد کلاینت و ۱۰۰ تکرار را تعریف می‌کند؛ فایل خلاصه، شمارش مستقل تکرارهای خام یا مشخصات کامل سخت‌افزار را ثبت نمی‌کند. بنابراین برای بازسازی دقیق زمان‌ها باید فرادادۀ اجرای جدید نیز ذخیره شود.

\begin{table}[htbp]
\centering
\small
\caption{میانگین و انحراف معیار زمان کل در فایل خلاصۀ ذخیره‌شده؛ واحد میلی‌ثانیه}
\label{tab:sbhe-total}
\begin{tabular}{rrrr}
\toprule
بیت & کلاینت & میانگین کل & انحراف معیار \\
\midrule
256 & 100 & 119.0243 & 31.5026 \\
256 & 200 & 188.0846 & 12.0490 \\
256 & 300 & 275.6820 & 11.6885 \\
512 & 100 & 152.4137 & 30.1520 \\
512 & 200 & 242.6216 & 17.3668 \\
512 & 300 & 338.7937 & 20.4376 \\
1024 & 100 & 411.2403 & 101.8251 \\
1024 & 200 & 521.9208 & 85.5255 \\
1024 & 300 & 723.1441 & 135.7782 \\
\bottomrule
\end{tabular}
\end{table}

\begin{table}[htbp]
\centering
\small
\caption{میانگین زمان مرحله‌های پروتکل در اجرای ذخیره‌شده؛ واحد میلی‌ثانیه}
\label{tab:sbhe-phases}
\resizebox{\textwidth}{!}{%
\begin{tabular}{rrrrrrr}
\toprule
بیت & کلاینت & آماده‌سازی & ثبت هویت & رمزگذاری & تجمیع/بررسی & رمزگشایی \\
\midrule
256 & 100 & 33.8433 & 1.6618 & 27.2437 & 55.1685 & 1.0871 \\
256 & 200 & 43.6925 & 2.7480 & 45.9302 & 94.1718 & 1.5322 \\
256 & 300 & 58.9859 & 4.2307 & 68.6870 & 141.5319 & 2.2465 \\
512 & 100 & 66.0327 & 1.5657 & 31.4395 & 50.5998 & 2.7661 \\
512 & 200 & 77.9172 & 2.8062 & 60.8421 & 96.9996 & 4.0564 \\
512 & 300 & 92.4105 & 4.1395 & 91.9830 & 144.9288 & 5.3018 \\
1024 & 100 & 293.5955 & 1.4030 & 53.0456 & 52.4542 & 10.7421 \\
1024 & 200 & 292.6545 & 2.8529 & 107.3289 & 104.4929 & 14.5817 \\
1024 & 300 & 352.8930 & 4.5927 & 177.0415 & 168.3493 & 20.2576 \\
\bottomrule
\end{tabular}}
\end{table}


در اندازۀ ۲۵۶ بیت، میانگین زمان کل برای ۱۰۰، ۲۰۰ و ۳۰۰ کلاینت به‌ترتیب ۱۱۹٫۰۲۴۳، ۱۸۸٫۰۸۴۶ و ۲۷۵٫۶۸۲۰ میلی‌ثانیه است. رشد تعداد کلاینت از ۱۰۰ به ۳۰۰، هزینۀ رمزگذاری و اعتبارسنجی/تجمیع را افزایش می‌دهد. در ۵۱۲ بیت، زمان‌های متناظر ۱۵۲٫۴۱۳۷، ۲۴۲٫۶۲۱۶ و ۳۳۸٫۷۹۳۷ میلی‌ثانیه و در ۱۰۲۴ بیت، ۴۱۱٫۲۴۰۳، ۵۲۱٫۹۲۰۸ و ۷۲۳٫۱۴۴۱ میلی‌ثانیه ثبت شده‌اند. این اعداد به پیکربندی‌های همین فایل محدودند و زمان پاسخ شبکه، آموزش مدل یا استنتاج متن را شامل نمی‌شوند.

\begin{figure}[htbp]
\centering\includegraphics[width=\textwidth]{original/sbhe-performance.png}
\caption{مقایسۀ عملکرد زمانی \en{SBHE-Benaloh} در اندازه‌های کلید و تعداد کلاینت‌های اجرای ذخیره‌شده}
\label{fig:sbhe-performance}
\end{figure}

برای ۱۰۰ کلاینت در ۲۵۶ بیت، میانگین زمان آماده‌سازی ۳۳٫۸۴۳۳، ثبت هویت ۱٫۶۶۱۸، رمزگذاری ۲۷٫۲۴۳۷، تجمیع و بررسی محموله‌ها ۵۵٫۱۶۸۵ و رمزگشایی ۱٫۰۸۷۱ میلی‌ثانیه است. اختلاف جزئی جمع زمان مرحله‌ها با زمان کل می‌تواند از سربار بخش‌های اندازه‌گیری‌نشده و گرد کردن ناشی شود؛ نباید با تغییر اعداد دادۀ اصلی پنهان شود. سهم قابل‌توجه تجمیع/بررسی نشان می‌دهد هزینۀ پروتکل فقط ضرب چند عدد نیست.

\begin{figure}[htbp]
\centering\includegraphics[width=\textwidth]{original/sbhe-phases.png}
\caption{مقایسۀ زمان مرحله‌های پروتکل در پیکربندی‌های ثبت‌شده}
\label{fig:sbhe-phases}
\end{figure}
\begin{figure}[htbp]
\centering\includegraphics[width=\textwidth]{original/sbhe-phase-shares.png}
\caption{توزیع هزینۀ مرحله‌ها با تغییر اندازۀ کلید؛ واحد زمان در نمودارها مشخص شده است}
\label{fig:sbhe-phase-shares}
\end{figure}
\begin{figure}[htbp]
\centering\includegraphics[width=\textwidth]{original/sbhe-scalability.png}
\caption{اثر تعداد کلاینت و طول کلید بر هزینۀ اجرای ذخیره‌شدۀ پروتکل}
\label{fig:sbhe-scalability}
\end{figure}

افزایش طول کلید، هزینۀ توان‌رسانی‌های پیمانه‌ای را افزایش می‌دهد. تعداد کلاینت نیز تعداد رمزگذاری‌ها و بررسی امضاها را بیشتر می‌کند؛ در مقابل، رمزگشاییِ نهایی یک تجمیع است. با این حال، شکل رشد باید از دادۀ چند اجرای کنترل‌شده بررسی شود؛ سه نقطۀ اندازۀ کلید برای نسبت دادن قانون پیچیدگی دقیق کافی نیست. همچنین طول کلید و ظرفیت پیام دو محور مستقل‌اند و کوچک بودن زمان، سطح امنیت مناسب را اثبات نمی‌کند.

\section{آزمون‌های کارکردی و مسیرهای توسعه}
مجموعۀ آزمون موجود، رمزگذاری و رمزگشایی، مجموع وزن‌دار، میانگین احرازاصالت‌شده، رد عضو غیرمجاز، رد تغییر وزن و متن رمز، رد بازپخش، رد دور نادرست و کنترل دقت اعشاری را بررسی می‌کند. هر نه آزمون در آماده‌سازی محیط اجرا و موفق شدند. این نتیجه، شاهد مشخص کارکردهای آزموده‌شده است؛ اثبات رسمی امنیت، آزمون استقرار شبکه یا پوشش همۀ ورودی‌های مخرب نیست.

برای توسعه، رمزگشایی آستانه‌ای می‌تواند اعتماد به مالک منفرد کلید را کاهش دهد؛ چنین تغییری به طراحی و تحلیل جداگانه نیاز دارد. رجیستری مبتنی بر گواهی، وضعیت پایدار دور و \en{nonce}، سیاست وزنِ قابل‌راستی‌آزمایی، اثبات بازه و سنجش هزینۀ پیام بزرگ‌تر از نیازهای اتصال به سناریوی مشارکتی‌اند. امضاهای تجمیع‌پذیر مانند \en{BLS} نیز یک گزینۀ پژوهشی‌اند، اما انتخاب آن‌ها بدون مدل تهدید، مقابله با کلید نامعتبر و سنجش هزینۀ تأیید، بهبود رایگان و قطعی محسوب نمی‌شود.

\section{جمع‌بندی}
دستاورد این فصل، صورت‌بندی و پیاده‌سازی قابل‌آزمونِ احراز اصالت محموله پیش از تجمیع همریخت است. درستی میانگین به اعتبار ورودی و نبود سرریز وابسته است؛ احراز اصالت، صحت معنایی پیش‌بینی را ثابت نمی‌کند؛ و مالکیت کلید خصوصی، مرز محرمانگی در برابر سرور را تعیین می‌کند. نتایج عددی نشان می‌دهند این ساختار در تنظیم‌های پژوهشیِ ثبت‌شده قابل اجراست. تبدیل آن به سامانۀ بالینی یا سرویس امنِ تولیدی، به ارزیابی و توسعۀ بیشتری نیاز دارد.
\FloatBarrier
